\chapter{Strong-Verb Stem Changes}

\section{What do regular, irregular, weak, and strong mean?}

The terms describe two related ideas. \textbf{Weak} and \textbf{strong}
describe how a verb forms its past tenses, whereas \textbf{regular} and
\textbf{irregular} describe how predictable its forms are. In this book,
\textit{regular} normally means a weak verb that follows the expected pattern.

\begin{center}
\small
\rowcolors{2}{referenceBlueBg}{white}
\begin{tabularx}{\textwidth}{@{}>{\bfseries}p{0.17\textwidth} X p{0.31\textwidth}@{}}
\toprule
Label & Meaning & Example forms \\
\midrule
Weak / regular & Adds \textit{-te} in the Präteritum and normally
ends the Partizip II in \textit{-t}; there is no unpredictable stem-vowel
change. & \textit{machen -- machte -- gemacht} \\
Strong & Changes its stem vowel in the Präteritum and normally ends the
Partizip II in \textit{-en}. Its key forms must be learnt. &
\textit{sprechen -- sprach -- gesprochen} \\
Mixed & Combines a stem change with the weak endings \textit{-te} and
\textit{-t}. & \textit{denken -- dachte -- gedacht} \\
Irregular & An umbrella label for a verb with forms that cannot all be
predicted from the normal weak pattern. Strong and mixed verbs are therefore
irregular, as are special verbs such as \textit{sein}. &
\textit{sein -- war -- gewesen} \\
\bottomrule
\end{tabularx}
\end{center}

\begin{tcolorbox}[
  colback=referenceYellowBg,
  colframe=genderPlural,
  title={How to read the labels in the verb tables},
  fonttitle=\bfseries
]
\textbf{Regular and weak usually mean the same thing in this book.}
\textbf{Strong verbs are irregular}, but not every irregular verb is strong:
mixed verbs and highly irregular verbs form other groups. A label such as
\textit{reflexive}, \textit{separable}, or \textit{+ Dativ} gives different
information and does not determine whether the verb is weak or strong.

Some regular patterns have predictable details: verbs ending in
\textit{-ieren} form the Partizip II without \textit{ge-}
(\textit{studiert}), and separable verbs place \textit{ge-} after the prefix
(\textit{aufgemacht}).
\end{tcolorbox}

Strong verbs form the Präteritum by changing the stem vowel rather than by
adding \textit{-te}. Their Partizip II normally ends in \textit{-en}. Some also
change their vowel in the second- and third-person singular Präsens. Because
the changes are not completely predictable, learn each verb through its key
forms:

\begin{tcolorbox}[
  colback=referenceBlueBg,
  colframe=genderMasculine,
  title={Recommended learning pattern},
  fonttitle=\bfseries
]
\centering
\textbf{Infinitiv} $\longrightarrow$ \textbf{er/sie/es im Präsens}
$\longrightarrow$ \textbf{Präteritum} $\longrightarrow$ \textbf{Perfekt}

\medskip
\textit{sprechen $\longrightarrow$ spricht $\longrightarrow$ sprach
$\longrightarrow$ hat gesprochen}
\end{tcolorbox}

\section{Learning shortcuts}

\begin{itemize}[leftmargin=*]
  \item Learn the third-person singular Präsens whenever it changes:
        \textit{nehmen -- er nimmt}, not merely \textit{nehmen}.
  \item Learn the Präteritum and Perfekt together:
        \textit{sprechen -- sprach -- hat gesprochen}.
  \item Mark verbs that normally use \textit{sein}, especially movement and
        change-of-state verbs such as \textit{gehen, kommen, bleiben}, and
        \textit{fallen}.
  \item Learn a prefixed verb with its family, but check whether the prefix is
        separable: \textit{aufstehen -- stand auf -- ist aufgestanden}, but
        \textit{verstehen -- verstand -- hat verstanden}.
\end{itemize}

\section{Strong verbs with a Präsens stem change}

The vowel change normally appears only with \textit{du} and
\textit{er/sie/es}. The plural forms retain the infinitive vowel:
\textit{ich fahre, du fährst, er fährt, wir fahren}.

\begin{center}
\small
\rowcolors{2}{referenceBlueBg}{white}
\begin{tabularx}{\textwidth}{@{}>{\bfseries}l l l l X@{}}
\toprule
Infinitive & English & er/sie/es & Präteritum & Perfekt \\
\midrule
empfehlen & recommend & empfiehlt & empfahl & hat empfohlen \\
essen & eat & isst & aß & hat gegessen \\
fahren & drive/travel & fährt & fuhr & ist/hat gefahren \\
fallen & fall & fällt & fiel & ist gefallen \\
fangen & catch & fängt & fing & hat gefangen \\
fressen & eat (animals) & frisst & fraß & hat gefressen \\
geben & give & gibt & gab & hat gegeben \\
gefallen & please & gefällt & gefiel & hat gefallen \\
gelten & be valid & gilt & galt & hat gegolten \\
halten & hold/stop & hält & hielt & hat gehalten \\
helfen & help & hilft & half & hat geholfen \\
laden & load/invite & lädt & lud & hat geladen \\
lassen & let/leave & lässt & ließ & hat gelassen \\
laufen & run/walk & läuft & lief & ist gelaufen \\
lesen & read & liest & las & hat gelesen \\
nehmen & take & nimmt & nahm & hat genommen \\
raten & advise/guess & rät & riet & hat geraten \\
schlafen & sleep & schläft & schlief & hat geschlafen \\
schlagen & hit & schlägt & schlug & hat geschlagen \\
sehen & see & sieht & sah & hat gesehen \\
sprechen & speak & spricht & sprach & hat gesprochen \\
stechen & sting/stab & sticht & stach & hat gestochen \\
tragen & carry/wear & trägt & trug & hat getragen \\
treffen & meet/hit & trifft & traf & hat getroffen \\
treten & step/kick & tritt & trat & ist/hat getreten \\
vergessen & forget & vergisst & vergaß & hat vergessen \\
waschen & wash & wäscht & wusch & hat gewaschen \\
werfen & throw & wirft & warf & hat geworfen \\
werden & become & wird & wurde & ist geworden \\
\bottomrule
\end{tabularx}
\end{center}

\section{Strong verbs without a Präsens stem change}

These verbs look regular in the Präsens but change their stem in the
Präteritum, the Partizip II, or both.

\begin{center}
\small
\rowcolors{2}{referenceGreenBg}{white}
\begin{tabularx}{\textwidth}{@{}>{\bfseries}l l l l X@{}}
\toprule
Infinitive & English & er/sie/es & Präteritum & Perfekt \\
\midrule
beginnen & begin & beginnt & begann & hat begonnen \\
beißen & bite & beißt & biss & hat gebissen \\
bieten & offer & bietet & bot & hat geboten \\
bitten & ask/request & bittet & bat & hat gebeten \\
bleiben & remain & bleibt & blieb & ist geblieben \\
finden & find & findet & fand & hat gefunden \\
fliegen & fly & fliegt & flog & ist/hat geflogen \\
fliehen & flee & flieht & floh & ist geflohen \\
gehen & go & geht & ging & ist gegangen \\
gießen & pour & gießt & goss & hat gegossen \\
genießen & enjoy & genießt & genoss & hat genossen \\
gewinnen & win & gewinnt & gewann & hat gewonnen \\
heben & lift & hebt & hob & hat gehoben \\
heißen & be called & heißt & hieß & hat geheißen \\
kommen & come & kommt & kam & ist gekommen \\
leihen & lend/borrow & leiht & lieh & hat geliehen \\
rufen & call & ruft & rief & hat gerufen \\
scheinen & shine/seem & scheint & schien & hat geschienen \\
schließen & close & schließt & schloss & hat geschlossen \\
schreiben & write & schreibt & schrieb & hat geschrieben \\
sitzen & sit & sitzt & saß & hat/ist gesessen \\
stehen & stand & steht & stand & hat/ist gestanden \\
steigen & climb/rise & steigt & stieg & ist gestiegen \\
streichen & paint/cancel & streicht & strich & hat gestrichen \\
trinken & drink & trinkt & trank & hat getrunken \\
tun & do & tut & tat & hat getan \\
verzeihen & forgive & verzeiht & verzieh & hat verziehen \\
ziehen & pull/move & zieht & zog & ist/hat gezogen \\
\bottomrule
\end{tabularx}
\end{center}

\section{Common prefixed families}

Many prefixed verbs preserve the stem change of their base verb. Inseparable
prefixes such as \textit{be-, emp-, ent-, er-, ver-}, and \textit{zer-} do not
add \textit{ge-} in the Partizip II.

\begin{center}
\small
\rowcolors{2}{referenceYellowBg}{white}
\begin{tabularx}{\textwidth}{@{}>{\bfseries}l l X@{}}
\toprule
Family & Pattern & Examples \\
\midrule
kommen & kam -- gekommen & bekommen: bekam -- bekommen \\
lassen & ließ -- gelassen & entlassen: entließ -- entlassen \\
leihen & lieh -- geliehen & verleihen: verlieh -- verliehen \\
nehmen & nahm -- genommen & übernehmen: übernahm -- übernommen; unternehmen: unternahm -- unternommen \\
schreiben & schrieb -- geschrieben & beschreiben: beschrieb -- beschrieben; unterschreiben: unterschrieb -- unterschrieben \\
sprechen & sprach -- gesprochen & widersprechen: widersprach -- widersprochen \\
stehen & stand -- gestanden & bestehen: bestand -- bestanden; entstehen: entstand -- entstanden; verstehen: verstand -- verstanden \\
halten & hielt -- gehalten & erhalten: erhielt -- erhalten; unterhalten: unterhielt -- unterhalten \\
ziehen & zog -- gezogen & beziehen: bezog -- bezogen \\
\bottomrule
\end{tabularx}
\end{center}

\section{Mixed verbs and modal verbs}

Mixed verbs combine a stem change with the weak \textit{-te} Präteritum and
the \textit{-t} Partizip ending. Modal verbs follow a closely related irregular
pattern and lose their umlaut in the Präteritum.

\begin{center}
\small
\rowcolors{2}{referenceRedBg}{white}
\begin{tabularx}{\textwidth}{@{}>{\bfseries}l l l l X@{}}
\toprule
Infinitive & English & er/sie/es & Präteritum & Perfekt \\
\midrule
bringen & bring & bringt & brachte & hat gebracht \\
denken & think & denkt & dachte & hat gedacht \\
kennen & know/be familiar with & kennt & kannte & hat gekannt \\
nennen & name/call & nennt & nannte & hat genannt \\
wissen & know a fact & weiß & wusste & hat gewusst \\
\addlinespace
dürfen & be allowed to & darf & durfte & hat gedurft \\
können & be able to & kann & konnte & hat gekonnt \\
mögen & like & mag & mochte & hat gemocht \\
müssen & have to & muss & musste & hat gemusst \\
sollen & should/be supposed to & soll & sollte & hat gesollt \\
wollen & want & will & wollte & hat gewollt \\
\bottomrule
\end{tabularx}
\end{center}
